\section{4.7.Echidna adaptar}\label{echidna-adaptar}

\subsection{4.7. Adaptador para Echidna y generación automática de
propiedades}\label{adaptador-para-echidna-y-generaciuxf3n-automuxe1tica-de-propiedades}

Una vez finalizado el proceso de correlación, EVMAudit incorpora una
fase adicional orientada a la validación dinámica de determinadas
vulnerabilidades mediante fuzzing.

Para ello, se desarrolló un adaptador específico encargado de
transformar los hallazgos correlacionados en propiedades compatibles con
Echidna, permitiendo complementar los resultados obtenidos mediante
análisis estático y ejecución simbólica.

La incorporación de esta fase responde a uno de los objetivos
específicos planteados en el trabajo, consistente en combinar distintas
técnicas de análisis con el fin de obtener una visión más completa del
estado de seguridad del contrato analizado.

!TODO: referencia cruzada a la Sección 3.2 (objetivos específicos).

\subsubsection{4.7.1. Objetivos del
adaptador}\label{objetivos-del-adaptador}

El módulo desarrollado persigue los siguientes objetivos:

\begin{itemize}
\tightlist
\item
  reutilizar la información obtenida durante la fase de correlación
\item
  generar automáticamente propiedades para Echidna
\item
  reducir la intervención manual del auditor
\item
  proporcionar una fase adicional de validación
\item
  mantener desacopladas las distintas fases del análisis.
\end{itemize}

La existencia de este componente permite integrar el fuzzing dentro del
pipeline general de EVMAudit sin introducir dependencias directas entre
las herramientas utilizadas.

\subsubsection{4.7.2. Generación de
wrappers}\label{generaciuxf3n-de-wrappers}

Echidna requiere que las propiedades de seguridad estén implementadas
como funciones Solidity dentro de un contrato específico.

Con el objetivo de automatizar este proceso, EVMAudit genera
dinámicamente un contrato wrapper que hereda del contrato original e
incorpora las propiedades correspondientes a las vulnerabilidades
detectadas.

El proceso seguido puede resumirse en las siguientes etapas:

\begin{enumerate}
\def\labelenumi{\arabic{enumi}.}
\tightlist
\item
  Recepción de los hallazgos correlacionados.
\item
  Consulta del catálogo SWC.
\item
  Obtención de las plantillas asociadas.
\item
  Sustitución de parámetros.
\item
  Generación del wrapper final.
\item
  Ejecución de Echidna.
\end{enumerate}

Gracias a esta aproximación, la generación de pruebas se realiza de
forma completamente automática.

!TODO: insertar diagrama del proceso de generación de wrappers.

!TODO: insertar referencia a la Figura X.

\subsubsection{4.7.3. Plantillas de
propiedades}\label{plantillas-de-propiedades}

Cada vulnerabilidad soportada dispone de una plantilla específica que
describe la propiedad que deberá ser evaluada por Echidna.

Estas plantillas permiten traducir información abstracta procedente del
análisis estático a estructuras ejecutables compatibles con el motor de
fuzzing.

La utilización de plantillas aporta varias ventajas:

\begin{itemize}
\tightlist
\item
  reutilización de código
\item
  independencia respecto a las herramientas de detección
\item
  facilidad de mantenimiento
\item
  incorporación sencilla de nuevas vulnerabilidades.
\end{itemize}

!TODO: insertar ejemplo de una plantilla real.

\subsubsection{4.7.4. Vulnerabilidades testables y no
testables}\label{vulnerabilidades-testables-y-no-testables}

Durante el desarrollo se observó que no todas las vulnerabilidades
pueden verificarse automáticamente mediante fuzzing.

Por este motivo, el adaptador distingue entre:

\paragraph{Vulnerabilidades testables}\label{vulnerabilidades-testables}

Son aquellas que pueden expresarse mediante propiedades directamente
evaluables por Echidna.

Algunos ejemplos incluyen:

\begin{itemize}
\tightlist
\item
  determinadas restricciones de acceso
\item
  invariantes económicas
\item
  comprobaciones relacionadas con balances
\item
  validaciones de lógica de negocio.
\end{itemize}

\paragraph{Vulnerabilidades no
testables}\label{vulnerabilidades-no-testables}

Corresponden a aquellas situaciones cuya explotación requiere
condiciones externas o escenarios complejos difíciles de reproducir
automáticamente.

Entre ellas destacan:

\begin{itemize}
\tightlist
\item
  ataques de reentrancia
\item
  determinados problemas de interacción entre contratos
\item
  escenarios dependientes del contexto de despliegue.
\end{itemize}

En estos casos, EVMAudit conserva la información asociada a la
vulnerabilidad y genera advertencias específicas, pero no intenta forzar
una validación automática que pudiera producir resultados incorrectos.

Esta decisión se adoptó con el objetivo de evitar conclusiones erróneas
y mantener un comportamiento conservador.

\subsubsection{4.7.5. Integración con
Echidna}\label{integraciuxf3n-con-echidna}

Una vez generado el wrapper correspondiente, el módulo Runner ejecuta
Echidna y recupera los resultados obtenidos durante el proceso de
fuzzing.

La salida generada por la herramienta se procesa posteriormente para
incorporarla al informe final.

De esta manera, EVMAudit consigue integrar tres técnicas de análisis
diferentes:

\begin{itemize}
\tightlist
\item
  análisis estático
\item
  ejecución simbólica
\item
  fuzzing.
\end{itemize}

La combinación de estas aproximaciones permite aprovechar las fortalezas
de cada una de ellas y compensar parcialmente sus limitaciones
individuales.

!TODO: referencia cruzada a la Sección 2.X (análisis estático).

!TODO: referencia cruzada a la Sección 2.X (ejecución simbólica).

!TODO: referencia cruzada a la Sección 2.X (fuzzing).

\subsubsection{4.7.6. Beneficios de la generación
automática}\label{beneficios-de-la-generaciuxf3n-automuxe1tica}

La generación automática de propiedades aporta diversas ventajas:

\begin{itemize}
\tightlist
\item
  reducción del trabajo manual
\item
  mayor reutilización de resultados
\item
  integración transparente con Echidna
\item
  facilidad para incorporar nuevas vulnerabilidades
\item
  mejora del flujo general de análisis.
\end{itemize}

No obstante, la validación automática obtenida no sustituye la revisión
manual realizada por el auditor, sino que constituye una capa adicional
de apoyo.
